\scanpage{16}
\hfill Rédaction n° 398
\begin{center}
COMPLÉMENTS AU CHAPITRE VII D'ALGÈBRE

DUALITÉ DES MODULES DE LONGUEUR FINIE SUR UN ANNEAU PRINCIPAL

-:-:-:-

CONTRE-RÉDACTION
\end{center}

[On se refuse à parler de modules injectifs, et d'autre part on se borne aux modules de longueur finie. Cela donne une rédaction minimale, (dans l'ensemble des rédactions $\neq\varnothing$)].

\uline{Notations}.

$A$ désigne un anneau principal, $K$ son corps des fractions. On suppose $K\neq A$ : i.e., $A$ possède au moins un élément extrémal.

Si $E$ est un $A$-module, on pose $E^*=\operatorname{Hom}(E,K/A)$ ; c'est un $A$-module. A tout homomorphisme $E\longrightarrow F$ correspond un homomorphisme $F^*\longrightarrow E^*$. Si $x\in E$, $x^*\in E^*$, on pose $x^*(x)=(x,x^*)$. On définit ainsi une application $A$-bilinéaire de $E\times E^*$ dans $K/A$, d'où un homomorphisme :
\[i_E:E\longrightarrow E^{**}.\]

PROPOSITION 1.- \uline{Si} $E$ \uline{est un} $A$-\uline{module de longueur finie}, $E^*$ \uline{est isomorphe}, (non canoniquement en général), \uline{à} $E$, \uline{et l'application} $i_E:E\longrightarrow E^{**}$ \uline{est un isomorphisme}.

Vu les théorèmes de structure, on peut se borner au cas où $E=A/tA$, avec $t\neq0$. Dans ce cas $E^*$ s'identifie au sous-module de $K/A$ formé des $x$ tels que $tx=0$, i.e. au sous-module $t^{-1}A/A$, isomorphe à $A/tA$. D'où l'isomorphie de $E$ et de $E^*$. Enfin, lorsque $E=A/tA$, la détermination précédente montre que $E^{**}$ s'identifie à $A/tA$ et l'on vérifie tout de suite que, dans ce cas, $i_E$ est l'application identique. D'où la deuxième partie de la proposition.
\mspage{15}
